近年, メディアコンテンツの流動性が急速に拡大し, メディアコンテンツ需要が高まっている. これに伴いストーリー創作の困難が増加し, ストーリーの創作支援が求められるている. 大規模言語モデルをはじめとしたAI技術の発展が著しく, これらが創作活動の支援に用いられると期待されているが, その活用方法は未だ確立されていない. 特にAIに不慣れな人でも有用な支援を受けられるようにするシステムの構築が求められている. そこで我々は, 物語構造分析に基づく LLM の制御を, 直感的な操作によって実現するインターフェースと, プロットの質向上を目的としたLLM とのインタラクティブな共創手法を提供する web アプリケーションを開発した. 本アプリケーションは, 手塚治虫原作の漫画『ブラック・ジャック』の物語構造分析に基づいたプロット生成を行う機能と, LLM とのインタラクティブなやり取りによってプロットの質を向上させることができる機能を備えており, LLM との共創手法を提供する. 本アプリケーションをプロのストーリークリエーターに利用してもらった結果, 操作性や指示の反映において高い評価を受け, また, インタラクティブな共創手法として有用であることが示唆された. 我々は, 本研究により,ストーリー創作支援における人とAIとの共創のあり方を探求した.

In recent years, the fluidity of media content has expanded rapidly, leading to a growing demand for media content. As a result, the challenges in creating stories have increased, and there is a growing need for story creation support. The development of AI technology, including large language models (LLMs), has been remarkable, and there are high expectations for these tools to aid in creative activities, though their utilization methods are not yet established. In particular, there is a demand for systems that offer meaningful support to individuals who are not familiar with AI. To address this, we have developed a web application that provides an intuitive interface for controlling LLMs based on narrative structure analysis, as well as interactive co-creation methods aimed at improving the quality of plot development. Our application includes features for generating plots based on the narrative structure analysis of Osamu Tezuka's manga "Black Jack" and functions to enhance plot quality through interactive interactions with the LLM, thus offering a collaborative co-creation approach. When utilized by professional story creators, our application received high evaluations for its usability and reflection of instructions, and it has been suggested as a useful interactive co-creation method. Through our study, we have explored approaches to co-creation between people and AI in support of story creation.

DeepLによる再翻訳
近年、メディアコンテンツの流動性が急速に拡大し、メディアコンテンツへの需要が高まっている。それに伴い、ストーリー作成の課題も増え、ストーリー作成支援のニーズが高まっている。大規模言語モデル（LLM）をはじめとするAI技術の発展は目覚ましく、その活用方法は確立されていないものの、創作活動を支援するツールとしての期待は大きい。特に、AIに不慣れな個人にとって有意義な支援を提供するシステムが求められている。そこで我々は、物語構造解析に基づくLLMを直感的に制御するためのインターフェースと、プロット作成の質の向上を目指した対話的な共創手法を提供するWebアプリケーションを開発した。本アプリケーションは、手塚治虫の漫画「ブラック・ジャック」の物語構造解析に基づくプロット生成機能と、LLMとの対話的なインタラクションによりプロットの質を高める機能を備え、協調的な共創アプローチを提供します。本アプリケーションをプロのストーリークリエイターに活用してもらったところ、その使い勝手や指示の反映性などが高く評価され、インタラクティブな共創手法として有用であることが示唆された。本研究を通して、物語創作を支援する人とAIの共創のアプローチを探ることができた。